\documentclass[12pt]{article}

\usepackage[utf8]{inputenc}
\usepackage[T1]{fontenc}

\title{Systematics}
\author{Luis Gonzalez}
\date{\today}

\begin{document}

\maketitle

\section{Introducción}
La familia Labridae es una de las familias más diversas y especializadas
de los grupos peces marinos asociados a los ambientes costeros. Algunos
de sus de estos peces presentan una marcada diversidad morfologica, fubncional
y ecologica siendo estas evidenciadas en la diferencia en tamaño corporal, 
forma del cuerpo, estructuras oseas, dentición, coloración y metodos 
de alimeentación. Esto ha permitido que los labridos esten presentes 
en una gran variedad de nichos, principalmente siendo los arrecifes de 
coral, además de esta gran variedad ha hecho que estos sean un grupo de interes 
para estudios de evolución y sistematica. Los primeros estudios filogeneticos de la 
familia se basaron principalmente de caracteres morfologicos mientras que posterirmente 
la incorporación de información molecular permitio evaluar las relaciones entre los 
principales linajes y contrastar las clasificaciones tradicionales.


Dentro de esta famlia esta 

La evolución como pilar

\end{document}